\subsection{Declaração do Problema}
\label{sec:problema}

O problema abordado neste trabalho é duplo, abrangendo tanto uma necessidade do usuário final quanto um desafio de engenharia de software.

\subsubsection{O Problema do Usuário: Sobrecarga de Informação}

Na era digital, indivíduos e organizações estão imersos em um volume exponencial de conteúdo de mídia não estruturado, como vídeos de longa duração, podcasts, palestras e documentos extensos. O consumo linear deste conteúdo é um processo inerentemente ineficiente, que demanda um investimento de tempo significativo para a extração de informações essenciais. Estudantes, pesquisadores e profissionais enfrentam uma barreira de produtividade, pois a necessidade de assimilar conhecimento rapidamente colide com a natureza demorada do consumo de mídia tradicional. A ausência de ferramentas eficazes para a síntese automática e contextualizada deste conteúdo representa um gargalo crítico no acesso à informação.

\subsubsection{O Problema Arquitetural: Sistemas Rígidos e Acoplados}

O desenvolvimento de sistemas para resolver o problema acima apresenta um desafio arquitetural significativo. Abordagens tradicionais frequentemente resultam em sistemas monolíticos e fortemente acoplados, com as seguintes deficiências:

\begin{itemize}
    \item \textbf{Falta de Extensibilidade:} Os sistemas são frequentemente projetados com um foco rígido em um único tipo de mídia (e.g., vídeo). A introdução de novos formatos de entrada (como arquivos de áudio ou documentos PDF) ou a adição de novas etapas ao pipeline de processamento (como análise de sentimento) exigem refatorações complexas e arriscadas, violando o Princípio Aberto/Fechado.
    
    \item \textbf{Acoplamento Tecnológico:} A lógica de negócio fica frequentemente acoplada a implementações de serviços específicos (e.g., um único provedor de nuvem para armazenamento ou um modelo de IA específico para transcrição). A substituição de uma dessas tecnologias por uma alternativa mais moderna ou de menor custo torna-se uma tarefa de alta complexidade.
    
    \item \textbf{Baixa Resiliência e Escalabilidade:} Em um fluxo de processamento síncrono ou rigidamente orquestrado, a falha em uma única etapa (como a indisponibilidade de uma API externa) pode comprometer todo o pipeline. Além disso, a incapacidade de escalar os estágios de processamento de forma independente (e.g., alocar mais recursos apenas para a etapa de transcrição, que é computacionalmente intensiva) limita a performance do sistema sob alta carga.
\end{itemize}

Portanto, o desafio de engenharia não é apenas "processar uma mídia", mas sim "construir uma plataforma que possa processar qualquer mídia, de forma confiável, escalável e que possa evoluir ao longo do tempo". A solução para o problema do usuário só é sustentável se a arquitetura subjacente for robusta e flexível.
